\subsection{The monoid structure on D(A)}

PROPOSITION 3. Let \(\mathfrak{a}, \mathfrak{a}', \mathfrak{b}, \mathfrak{b}'\) be elements of I(A). The relations \(\mathfrak{a} > \mathfrak{a}'\) and \(\mathfrak{b} > \mathfrak{b}'\) imply \(\mathfrak{ab} > \mathfrak{a}'\mathfrak{b}'\) .

We may restrict our attention to the case where \(\mathfrak{b} = \mathfrak{b}'\) . Then let \(Ax\) be a fractional principal ideal containing \(a'b\) ; for every non-zero element \(y\) of \(\mathfrak{b}\) , \(Ax \supset a'y\) and hence \(Axy^{-1} \supset a'\) , whence \(Axy^{-1} \supset a\) and \(Ax \supset ay\) . Varying \(y\) , it is seen that \(Ax \supset ab\) , whence \(ab > a'b\) .

It follows from Proposition 3 that multiplication on I(A) defines, by passing to the quotient, a law of composition on D(A) which is obviously associative and commutative. It is written additively so that we may write:

\[
\operatorname{div} (a b) = \operatorname{div} a + \operatorname{div} b,\tag{2}
\]

for a, b in I(A). Clearly div(1) is an identity element for this addition; this element is denoted by 0. Proposition 3 proves further that the order structure on D(A) is compatible with this addition (Algebra, Chapter VI, § 1, no. 1) and, more precisely (no. 1, Proposition 2 (ii)):

\[
\begin{array}{r l} \inf (\operatorname{div} a + \operatorname{div} b, \operatorname{div} a + \operatorname{div} c) & = \inf (\operatorname{div} (a b), \operatorname{div} (a c)) = \operatorname{div} (a b + a c) \\ & = \operatorname{div} (a (b + c)) = \operatorname{div} a + \operatorname{div} (b + c) = \operatorname{div} a + \inf (\operatorname{div} b, \operatorname{div} c). \end{array}
\]

For a fractional ideal \(a \neq 0\) to be such that div \(a \geqslant 0\) in D(A), it is necessary and sufficient that \(a \subset A\) (in other words, that a be an integral ideal of A).

For two elements \(x, y\) of \(\mathbf{K}^*\) , the relation \(\operatorname{div}(x) = \operatorname{div}(y)\) is equivalent to \(A x = A y\) ; the set of principal divisors of \(A\) with the order relation and the monoid law induced by that on \(D(A)\) is an ordered group canonically isomorphic to the multiplicative group of fractional principal ideals ordered by the opposite order relation to inclusion (Algebra, Chapter VI, § 1, no. 5). The relation \(S\) between two elements \(P, Q\) of \(D(A)\) :

\[
" \text { there   exists } x \in \mathbf {K} ^ {*} \text { such   that } \mathbf {P} = \mathbf {Q} + \operatorname{div} (x)"
\]

is therefore an equivalence relation since the relation \(P = Q + \text{div}(x)\) is equivalent to \(Q = P + \text{div}(x^{-1})\) ; if P and Q are congruent modulo S, they are called equivalent divisors of A. Clearly moreover the relation S is compatible with the law of the monoid D(A) and the latter therefore defines, by taking quotients, a monoid structure on D(A)/S; this monoid is called the divisor class monoid of A.

PROPOSITION 4. Let \(\mathfrak{a}\) , \(\mathfrak{b}\) be two divisorial fractional ideals of \(\mathbf{A}\) . The following properties are equivalent:

\begin{enumerate}
\def\labelenumi{(\alph{enumi})}
\item
  div a and div b are equivalent divisors;
\item
  there exists \(x \in \mathbf{K}^*\) such that \(\mathfrak{b} = xa\) .
\end{enumerate}

If \(\operatorname{div} \mathfrak{b} = \operatorname{div} \mathfrak{a} + \operatorname{div}(x)\) for some \(x \in \mathbf{K}^*\) , then \(\operatorname{div} \mathfrak{b} = \operatorname{div}(xa)\) and, as \(\mathfrak{b}\) and \(xa\) are divisorial, \(\mathfrak{b} = xa\) , which proves the proposition.

Let \(\mathfrak{a}\) be an invertible fractional ideal (Chapter II, § 5, no. 6); then \(\mathfrak{a} = \mathbf{A}\) : (A: \(\mathfrak{a}\) ) (loc. cit., Proposition 10) and hence \(\mathfrak{a}\) is divisorial (no. 1, Proposition 1). The group \(J(\mathbf{A})\) of invertible fractional ideals is therefore identified with a subgroup of the monoid \(D(\mathbf{A})\) and the canonical image of \(J(\mathbf{A})\) in \(D(\mathbf{A}) / S\) with the group of classes of projective A-modules of rank 1 (Chapter II, § 5, no. 7, Corollary 2 to Proposition 12 and Remark 1).

THEOREM 1. Let A be an integral domain. For the monoid D(A) of divisors of A to be a group, it is necessary and sufficient that A be completely integrally closed.

Suppose that D(A) is a group. Let \(x \in K\) ; suppose that A{[}x{]} is contained in a finitely generated sub-A-module of K. Then we have seen (no. 1) that \(a = A[x]\) is an element of I(A). Then \(xa \subset a\) and hence \(\text{div}(x) + \text{div}a \geqslant \text{div}a\) . Since D(A) is an ordered group, we conclude that \(\text{div}(x) \geqslant 0\) , whence \(x \in A\) . Thus A is completely integrally closed (Chapter V, § 1, no. 4, Definition 5).

Conversely, suppose that A is completely integrally closed. Let a be a divisorial ideal. We shall show that \(\operatorname{div} a + \operatorname{div}(A: a) = 0\) , which will prove that \(\mathrm{D}(A)\) is a group. As \(a(A: a) \subset A\) , it suffices (no. 1) to verify that every fractional principal ideal \(Ax^{-1}\) which contains \(a(A: a)\) also contains A. Now, for \(y \in K^{*}\) , the relation \(Ay \supset a\) implies \(y^{-1} \in A: a\) , whence \(y^{-1}a \subset a(A: a) \subset Ax^{-1}\) and hence \(xa \subset Ay\) . As a is divisorial, we deduce that \(xa \subset a\) , whence \(x^{n}a \subset a\) for all \(n \in N\) . There exist elements \(x_{0}, x_{1}\) of \(K^{*}\) such that \(Ax_{0} \subset a \subset Ax_{1}\) ; therefore \(x^{n}x_{0} \in Ax_{1}\) , whence \(x^{n} \in Ax_{1}x_{0}^{-1}\) . As A is completely integrally closed, \(x \in A\) , that is \(Ax^{-1} \supset A\) , which completes the proof.

Note that, if A is completely integrally closed (and even Noetherian), a divisorial ideal of A is not necessarily invertible, in other words, in general \(J(A) \neq D(A)\) (Exercise 2 and §3, no. 2, Proposition 1).

\pandocbounded{\includegraphics[keepaspectratio]{images/part003_63c67829.jpg}}

COROLLARY. Let A be a completely integrally closed domain and a divisorial fractional ideal of A. Then, for every fractional ideal \(\mathfrak{b} \neq 0\) of A, \(\operatorname{div}(\mathfrak{a}: \mathfrak{b}) = \operatorname{div} \mathfrak{a} - \operatorname{div} \mathfrak{b}\) .

By virtue of formula (1) of no. 1:

\[
\operatorname{div} (a: b) = \operatorname{div} \left(\bigcap_ {\gamma \in b, \gamma \neq 0} \gamma ^ {- 1} a\right) = \sup _ {\gamma \in b, \gamma \neq 0} \operatorname{div} (\gamma ^ {- 1} a)
\]

taking account of Proposition 2 and the fact that the fractional ideals \(y^{-1}\mathfrak{a}\) are divisorial. But since \(\mathrm{D(A)}\) is an ordered group (Algebra, Chapter VI, § 1, no. 8):

\[
\sup _ {\gamma \in b, \gamma \neq 0} \operatorname{div} (\gamma ^ {- 1} a) = \sup _ {\gamma \in b, \gamma \neq 0} (\operatorname{div} a - \operatorname{div} (\gamma)) = \operatorname{div} a - \inf _ {\gamma \in b, \gamma \neq 0} \operatorname{div} (\gamma) = \operatorname{div} a - \operatorname{div} b.
\]

